% SGA 3, Expose IV, English reconstruction.
% Scope: editor's note 38 and Lemmas 4.4.8.1--4.4.8.2, stopping before
% Proposition 4.4.9.
% Authority: Polo--Gille Exp4-13oct24.pdf, component-local page 41 partial
% (combined-reader page 229; printed page 217).
% The external English Markdown was used only as a comparison.

\setcounter{footnote}{37}

下面的引理将在 VI\textsubscript{B}, \SGARefLink{par:sga3-VIB-8.1}{8.1}
和 \SGARefLink{prop:sga3-VIB-8.2}{8.2} 中用到.
\footnote{编者注: 引理 \SGARefLink{sga3:iv:lemma:4:4:8:1}{4.4.8.1}
和 \SGARefLink{sga3:iv:lemma:4:4:8:2}{4.4.8.2} 是新增的.}

\medskip
\noindent\SGARefTarget{sga3:iv:lemma:4:4:8:1}\textbf{引理 4.4.8.1.}
\textit{设 \(Q\subset P\) 是 \(\mathscr C\) 上的群预层,
且 \(Q\) 在 \(P\) 中正规. 假设 \(P\) 分离而 \(Q\) 是层.
则群预层 \(P/Q\) 分离. 因此,}
\[
  a(P/Q)=L(P/Q)
  =\varinjlim_{R\in J(S)}(P/Q)(R).
\]

只须证明第一个断言, 因为第二个断言由它及
\SGARefLink{sga3:iv:proposition:4:3:11:item:iv:01}{4.3.11(iv)} 和
\SGARefLink{sga3:iv:proposition:4:3:11:item:ii:01}{(ii)} 得到.
设 \(S\in\operatorname{Ob}\mathscr C\), \(R\) 是 \(S\) 的筛,
\(y\in P(S)/Q(S)\) 在 \((P/Q)(R)\) 中的像为单位元.
须证 \(y=1\).

由假设, 态射 \(P(S)\to P(R)\) 是单射,
而 \(Q(S)\to Q(R)\) 是同构, 并且下列交换图的上行正合:
\[
\begin{array}{ccccccccc}
1&\longrightarrow&Q(S)&\longrightarrow&P(S)&\longrightarrow&P(S)/Q(S)
  &\longrightarrow&1\\
&&\big\downarrow{\scriptstyle\sim}&&\big\downarrow&&\big\downarrow\\
1&\longrightarrow&Q(R)&\longrightarrow&P(R)&\longrightarrow&(P/Q)(R).&&
\end{array}
\]
若下行正合, 就得到\SGARefLink{sga3:iv:lemma:4:4:8:1}{所求结果}.
设 \(f:R\to P\) 在 \((P/Q)(R)\) 中的像为单位元,
并取 \(\phi:T\to R\), 其中 \(T\in\operatorname{Ob}\mathscr C\).
则 \(f\circ\phi\) 在
\[
  (P/Q)(T)=P(T)/Q(T),
\]
中为单位元, 即 \(f\circ\phi\in Q(T)\).
因此 \(\phi\mapsto f\circ\phi\) 是函子性的映射 \(R(T)\to Q(T)\),
从而定义了一个函子态射 \(\pi:R\to Q\),
使 \(\pi(\operatorname{id}_R)=f\). 于是 \(f\in Q(R)\).
这证明了下行的正合性, 从而证明了\SGARefLink{sga3:iv:lemma:4:4:8:1}{此引理}.

\medskip
\noindent\SGARefTarget{sga3:iv:lemma:4:4:8:2}\textbf{引理 4.4.8.2.}
\textit{设 \(H\subset G\) 是 \(\mathscr C\) 上的群预层.}
\begin{enumerate}
  \item[(i)]\SGARefTarget{sga3:iv:lemma:4:4:8:2:item:i:01}
    \textit{若 \(H\) 在 \(G\) 中正规, 则 \(L(H)\) 在 \(L(G)\) 中正规.}
  \item[(ii)]\SGARefTarget{sga3:iv:lemma:4:4:8:2:item:ii:01}
    \textit{若 \(H\) 位于 \(G\) 的中心, 则 \(L(H)\) 位于 \(L(G)\) 的中心.}
\end{enumerate}

设 \(S\in\operatorname{Ob}\mathscr C\).
须证 (参见 I~\SGARefLink{sga3:i:definition:2:3:6}{2.3.6})
\(L(H)(S)\) 在 \(L(G)(S)\) 中正规, 分别地, 位于其中心. 取
\[
  h\in L(H)(S),\qquad g\in L(G)(S).
\]
存在 \(S\) 的一个筛 \(R\) 及元素 \(h'\in H(R)\),
\(g'\in G(R)\), 使得按 \SGARefLink{sga3:iv:lemma:4:3:10}{4.3.10} 的记号有
\[
  h=z_R(h'),\qquad g=z_R(g')
\]
由于 \(z_R\) 是群态射,
\[
  ghg^{-1}h^{-1}=z_R(g'h'g'^{-1}h'^{-1}).
\]
在情形 \SGARefLink{sga3:iv:lemma:4:4:8:2:item:i:01}{(i)} 中,
\[
  g'h'g'^{-1}h'^{-1}\in H(R),
\]
因而 \(ghg^{-1}h^{-1}\in L(H)(S)\).
在情形 \SGARefLink{sga3:iv:lemma:4:4:8:2:item:ii:01}{(ii)} 中,
\(g'h'g'^{-1}h'^{-1}=1\), 从而 \(ghg^{-1}h^{-1}=1\).

% Translation coverage: entire assigned source component.
% Source lines: 81; translated body lines before this audit comment: 86.
% Math preserved: 47 inline formulas, 7 display blocks.
% Reference targets: 4/4; links: 12/12; labels: 0/0; footnotes: 1/1.
% Unresolved TODOs: 0; compile and visual review delegated to the volume integrator.
